\phantomsection\label{sec:teaching-reflections}
\section*{教学思考}

\subsection*{教学理念}
我在教学中采用基于问题的学习模式，让学生面对实际问题，并引导他们形成结合理论与实践的学习方法。
我发现，当学生理解所学内容的重要性和实用价值时，学习效果最好。我将这些理念贯彻于课程建设和项目指导工作。

\subsection*{课程建设}
我为“AI Systems \& Infrastructure”课程以每个模块一篇博客的形式编写了完全原创的教学材料。
考虑到参与该课程的大多数学生缺少相关预备知识，我准备的材料以贴近学生背景的方式讲解相关概念。
博客形式便于持续更新，浏览方便且灵活，可以加入多媒体内容以适应学生多样的学习风格。

\vspace{\entryskip}
该课程的每个模块围绕现实世界的实际问题展开，并配有相互衔接的实践练习。
课程以小型项目收尾，培养学生综合运用课程所学全部技能的能力。
课堂教学力求富有启发性，并为学生后续的深入学习提供切入点。
口试重在考查学生能否解决实际问题，而非单纯记忆代码语法或复述概念名称。

\vspace{\entryskip}
在学校汇总学生反馈形成的学期评估报告中，该课程获得了良好的评价，学生着重肯定了课程材料的实用性。

\subsection*{项目指导}
我指导本科生和硕士生的学期项目，并撰写版式生动、引人入胜的项目提案，以实际问题为背景提供明确的项目方向。
我在项目指导中同样采用基于问题的学习模式。
我从实际问题和真实案例出发，着重说明规划项目需要考虑的关键因素，并提供充足的资源，让学生自行探索如何将这些因素落实到项目中。

\vspace{\entryskip}
我力求平易近人并为学生提供支持，同时为他们留出自主决策的空间。
我在学生遇到障碍时给予帮助，为技术路线提供建议，并鼓励他们主导项目方向。学生在独立应对项目挑战的过程中逐步建立信心，并形成适合自己的项目推进方式。
学生评价我的指导很有帮助，并认为我平易近人、易于沟通。

\subsection*{教学计划}
这些教学实践成效良好，我计划在此基础上持续改进。
我将继续遵循三项核心原则：扎实掌握所授内容，了解学生已有的经验和背景，以及贯彻基于问题的学习方法。同时，我将坚持已被证明有效的具体做法，包括博客式课程材料、富有启发性和参与感的课堂、公开透明的课程与考试形式和要求，以及循序渐进的模块安排。
我相信这套方法适用于各类计算机科学课程的教学。
我也有兴趣开设数据挖掘、人工智能等计算机相关方向的新课程。

\vspace{\entryskip}
我也已做好指导博士生的准备，并将沿用计算机科学和工程实践中普遍适用的原则，包括明确问题定义、批判性地分析设计选择、形成有说服力的结论，以及把控成果展示的质量。
我希望通过以科研为基础的教学，培养能够批判性地看待技术并运用知识解决实际问题的毕业生。
